\section{Introduction}
\label{sec:introduction}
Electroencephalography (EEG) provides deep insight into brain activity without requiring invasive procedures, and plays a crucial role in clinical diagnostics, cognitive neuroscience, and human-computer interaction. In recent years, deep neural networks have significantly advanced EEG analysis, shifting from handcrafted pipelines to end-to-end learning systems~\cite{craik2019deep}. Transformer-based models now rival traditional signal processing techniques by jointly modelling long-range temporal dynamics and cross-channel correlations~\cite{song2022eeg,wen2022transformers}. 

Despite this progress, \emph{a fundamental bottleneck remains}: EEG corpora exhibit significant \emph{topological heterogeneity}. Electrode count and placement vary widely across public and private datasets, making it difficult to transfer models across montages. This limitation manifests in pronounced performance degradation during cross-dataset evaluation. For example, motor-imagery decoders lose up to 14 percentage points (pp) in accuracy when transferring from PhysioNet to KU datasets~\cite{xu2020cross}, while state-of-the-art emotion-recognition models such as BIOT and MMM exhibit 13–15 pp drops between SEED and DEAP montages~\cite{yang2023biot,yi2023learning}. Similarly, patient-to-patient transfer in stereotactic EEG (sEEG) remains an unsolved challenge, with naive models performing near chance without explicit spatial encoding~\cite{mentzelopoulos2024neural}. 

Existing approaches offer limited solutions to this problem. Some train bespoke models for each montage, while others retain only shared electrodes---discarding up to \textbf{80\%} of available data~\cite{lin2023eeg}. More general approaches that flatten channels and time into long sequences incur quadratic self-attention complexity, \(\bigO\!\bigl((S \cdot C)^{2}\bigr)\) where $S$ is the number of time segments and $C$ is the number of electrodes (channels), rapidly exhausting memory on dense caps~\cite{yang2023biot}. These challenges underscore the need for a \textbf{single, montage-agnostic architecture that scales efficiently with electrode count}.

\textbf{LUNA} (\textbf{L}atent \textbf{U}nified \textbf{N}etwork \textbf{A}rchitecture) directly addresses this gap. Our key innovation is a topology-invariant encoder that maps arbitrary electrode layouts into a fixed latent space via learned queries and cross-attention. Temporal self-attention layers then operate exclusively on this latent space, decoupling computational cost from the number of electrodes. We pre-train LUNA using a masked-patch reconstruction objective on \textsc{TUEG} \cite{obeid2016temple} and \textsc{Siena} \cite{detti2020siena} (over \textit{21,000} hours of raw EEG data), and fine-tune on four downstream benchmarks spanning abnormality and artifact detection, slowing classification, and emotion recognition.

The key contributions of this work are the following:
\begin{itemize}[leftmargin=*]
        \item \textbf{Topology-invariant encoder.} A learnt query / cross-attention module that projects arbitrary-sized channel sets into a fixed latent space.
        \item \textbf{Linear-in-channels complexity.} Patch-wise temporal attention that decouples FLOPs and memory from electrode count.
        \item \textbf{State-of-the-art accuracy-efficiency trade-off.} LUNA achieves strong results across a range of EEG benchmarks, demonstrating significant capabilities with balanced accuracies of \textbf{81.57\%} on TUAB and \textbf{39.18\%} on SEED-V \cite{liu2021comparing}, and AUROC scores of \textbf{0.921} on TUAR and \textbf{0.802} on TUSL, while reducing FLOPs by \textbf{300$\times$} and GPU memory footprint by up to \textbf{10$\times$} on high-density EEG recordings. Crucially, these gains hold across diverse electrode configurations, confirming LUNA's generalization capability.
\end{itemize}

